\documentclass[aps,prb,onecolumn,nofootinbib]{revtex4-2}

\usepackage[T1]{fontenc}
\usepackage[utf8]{inputenc}

\usepackage{graphicx}
\usepackage{amsmath,amsfonts,amssymb,bm,mathtools}
\usepackage{algorithm}
\usepackage{algpseudocode}
\usepackage{float}
\usepackage{booktabs}
\usepackage{tabularx}
\usepackage{xcolor}

\makeatletter
\providecommand{\ALG@name}{Algorithm}
\makeatother

\usepackage[colorlinks=true,urlcolor=blue,linkcolor=blue,citecolor=blue]{hyperref}

\newcommand{\ket}[1]{|#1\rangle}
\newcommand{\bra}[1]{\langle #1|}
\newcommand{\braket}[2]{\langle #1 \mid #2 \rangle}
\newcommand{\Tr}{\mathrm{Tr}}
\newcommand{\Id}{\mathbb{I}}
\newcommand{\eps}{\epsilon}
\newcommand{\Order}{\mathcal{O}}

\begin{document}
	
	\title{\texorpdfstring{Numerical Implementation of Matrix Product States and DMRG\\
			\large Based on the review article ``The density-matrix renormalization group in the age of matrix product states'' by U.~Schollw\"ock, Annals of Physics \textbf{326}, 96--192 (2011)}
		{Numerical Implementation of Matrix Product States and DMRG}}
	
	\author{Kimi 2.5}
	\author{GPT 5.4}
	\date{\today}
	
	\maketitle
	
	\tableofcontents
	
	\section{Scope and conventions}
	
	These implementation notes summarize the practical coding structure of matrix product state (MPS), matrix product operator (MPO), and finite-system DMRG algorithms for one-dimensional systems with open boundary conditions (OBC). The presentation is based on Schollw\"ock's review article, but rewritten with explicit implementation conventions suitable for numerical work.
	
	We assume throughout:
	\begin{itemize}
		\item open boundary conditions,
		\item dense tensors without explicit quantum-number block sparsity,
		\item finite-system algorithms,
		\item NumPy-like row-major array storage.
	\end{itemize}
	
	We use the following tensor conventions throughout:
	\begin{itemize}
		\item MPS tensor on site $i$:
		\[
		A^{[i]}_{\alpha_{i-1},\sigma_i,\alpha_i},
		\]
		with shape $(D_{i-1}, d_i, D_i)$.
		\item MPO tensor on site $i$:
		\[
		W^{[i]}_{\beta_{i-1},\beta_i,\sigma_i,\sigma_i'},
		\]
		with shape $(\chi_{i-1}, \chi_i, d_i, d_i)$. We use the operator convention
		\[
		\hat W^{[i]}
		=
		\sum_{\sigma_i,\sigma_i'}
		W^{[i]}_{\beta_{i-1},\beta_i,\sigma_i,\sigma_i'}
		\ket{\sigma_i}\bra{\sigma_i'},
		\]
		so the first physical index is the ket index and the second is the bra index.
		\item OBC dimensions:
		\[
		D_0 = D_L = 1,\qquad \chi_0 = \chi_L = 1.
		\]
	\end{itemize}
	
	Unless stated otherwise, all tensors should be stored as \texttt{complex128}.
	
	\section{Reshape conventions}
	\label{sec:reshape}
	
	For numerical stability and reproducibility, reshape conventions must be fixed once and used everywhere.
	
	For a rank-3 MPS tensor
	\[
	T^{[i]}_{\alpha_{i-1},\sigma_i,\alpha_i},
	\]
	with storage shape $(D_{i-1},d_i,D_i)$, define the left-grouped matrix
	\begin{equation}
		T^{[i]}_{(\alpha_{i-1}\sigma_i),\alpha_i}
		\in \mathbb{C}^{(D_{i-1}d_i)\times D_i},
		\label{eq:left-group}
	\end{equation}
	and the right-grouped matrix
	\begin{equation}
		T^{[i]}_{\alpha_{i-1},(\sigma_i\alpha_i)}
		\in \mathbb{C}^{D_{i-1}\times(d_iD_i)}.
		\label{eq:right-group}
	\end{equation}
	
	In row-major (C-order) storage, the flattening used in Eq.~\eqref{eq:left-group} is
	\begin{equation}
		r = \alpha_{i-1} d_i + \sigma_i,\qquad c = \alpha_i.
	\end{equation}
	
	In NumPy:
	\begin{verbatim}
		M_left  = T.reshape((D_left * d, D_right), order='C')
		M_right = T.reshape((D_left, d * D_right), order='C')
	\end{verbatim}
	
	Typical inverse reshapes are
	\begin{verbatim}
		A  = U.reshape((D_left, d, D_new), order='C')
		B  = Vh.reshape((D_new, d, D_right), order='C')
	\end{verbatim}
	
	All QR/SVD steps must be consistent with these conventions, and the inverse reshape must use the same memory order and grouped-index convention.
	
	\section{Canonical forms}
	
	For fixed physical index $\sigma$, one may view $A^{[i]\sigma}$ as a matrix from the left bond to the right bond.
	
	Using index notation avoids ambiguity. A tensor $A^{[i]}$ is \emph{left-canonical} if
	\begin{equation}
		\sum_{\alpha_{i-1},\sigma_i}
		A^{[i]}_{\alpha_{i-1},\sigma_i,\alpha_i}
		A^{[i]*}_{\alpha_{i-1},\sigma_i,\alpha_i'}
		=
		\delta_{\alpha_i,\alpha_i'}.
		\label{eq:left-canonical}
	\end{equation}
	Equivalently, if $A^{[i]}$ is reshaped as in Eq.~\eqref{eq:left-group}, then
	\begin{equation}
		\left(A^{[i]}_{(\alpha_{i-1}\sigma_i),\alpha_i}\right)^\dagger
		A^{[i]}_{(\alpha_{i-1}\sigma_i),\alpha_i}
		= \Id_{D_i}.
	\end{equation}
	
	A tensor $B^{[i]}$ is \emph{right-canonical} if
	\begin{equation}
		\sum_{\sigma_i,\alpha_i}
		B^{[i]}_{\alpha_{i-1},\sigma_i,\alpha_i}
		B^{[i]*}_{\alpha_{i-1}',\sigma_i,\alpha_i}
		=
		\delta_{\alpha_{i-1},\alpha_{i-1}'}.
		\label{eq:right-canonical}
	\end{equation}
	Equivalently, if $B^{[i]}$ is reshaped as in Eq.~\eqref{eq:right-group}, then
	\begin{equation}
		B^{[i]}_{\alpha_{i-1},(\sigma_i\alpha_i)}
		\left(B^{[i]}_{\alpha_{i-1},(\sigma_i\alpha_i)}\right)^\dagger
		= \Id_{D_{i-1}}.
	\end{equation}
	
	A mixed-canonical state with orthogonality center at bond $i$ may be written as
	\begin{equation}
		\ket{\psi} =
		\sum_{\alpha_i}
		s^{[i]}_{\alpha_i}
		\ket{\alpha_i}_L \ket{\alpha_i}_R,
	\end{equation}
	where $s^{[i]}_{\alpha_i}\ge 0$ are Schmidt singular values and the left/right Schmidt states are orthonormal. For a normalized state,
	\[
	\sum_{\alpha_i} \left(s^{[i]}_{\alpha_i}\right)^2 = 1.
	\]
	
	\section{Canonicalization algorithms}
	
	\subsection{Left sweep by reduced QR}
	
	Given a general tensor $M^{[i]}$, reshape it as
	\[
	M^{[i]}_{(\alpha_{i-1}\sigma_i),\alpha_i}.
	\]
	Perform a reduced QR factorization
	\begin{equation}
		M^{[i]}_{(\alpha_{i-1}\sigma_i),\alpha_i} = Q R,
	\end{equation}
	where
	\[
	Q\in \mathbb{C}^{(D_{i-1}d_i)\times k},\qquad
	R\in \mathbb{C}^{k\times D_i},\qquad
	k=\min(D_{i-1}d_i,D_i).
	\]
	Reshape $Q$ into a left-canonical tensor
	\[
	A^{[i]}_{\alpha_{i-1},\sigma_i,\tilde\alpha_i},
	\]
	with new bond dimension $\tilde D_i = k$, then absorb $R$ into the next tensor:
	\begin{equation}
		M^{[i+1]}_{\tilde\alpha_i,\sigma_{i+1},\alpha_{i+1}}
		\leftarrow
		\sum_{\alpha_i}
		R_{\tilde\alpha_i,\alpha_i}
		M^{[i+1]}_{\alpha_i,\sigma_{i+1},\alpha_{i+1}}.
		\label{eq:left-qr-absorb}
	\end{equation}
	
	\begin{algorithm}[H]
		\caption{LeftCanonicalizeByQR}
		\begin{algorithmic}[1]
			\Require MPS tensors $\{M^{[1]},\dots,M^{[L]}\}$
			\Ensure Left-canonical tensors $\{A^{[1]},\dots,A^{[L-1]}\}$ and final center tensor $C^{[L]}$
			\For{$i=1$ to $L-1$}
			\State reshape $M^{[i]} \to \hat M$ with shape $(D_{i-1}d_i)\times D_i$
			\State compute reduced QR: $\hat M = Q R$
			\State reshape $Q \to A^{[i]}$ with shape $(D_{i-1},d_i,\tilde D_i)$
			\State absorb $R$ into $M^{[i+1]}$ via Eq.~\eqref{eq:left-qr-absorb}
			\EndFor
			\State set $C^{[L]} \gets M^{[L]}$
		\end{algorithmic}
	\end{algorithm}
	
	\paragraph{Important practical note.}
	For canonicalization, use \emph{reduced QR without pivoting}. Pivoted QR changes the bond basis by a permutation and is inconvenient unless the permutation is explicitly propagated into neighboring tensors.
	
	\subsection{Right sweep by reduced QR}
	
	To right-canonicalize site $i$, reshape
	\[
	M^{[i]}_{\alpha_{i-1},(\sigma_i\alpha_i)}
	\in \mathbb{C}^{D_{i-1}\times(d_iD_i)}.
	\]
	Compute a reduced QR factorization of the Hermitian transpose:
	\begin{equation}
		\left(M^{[i]}_{\alpha_{i-1},(\sigma_i\alpha_i)}\right)^\dagger = Q R,
	\end{equation}
	with
	\[
	Q\in\mathbb{C}^{(d_iD_i)\times k},\qquad
	R\in\mathbb{C}^{k\times D_{i-1}},\qquad
	k=\min(d_iD_i,D_{i-1}).
	\]
	Then
	\begin{equation}
		M^{[i]} = R^\dagger Q^\dagger.
	\end{equation}
	Reshape $Q^\dagger$ into a right-canonical tensor
	\[
	B^{[i]}_{\tilde\alpha_{i-1},\sigma_i,\alpha_i},
	\]
	with new bond dimension $\tilde D_{i-1}=k$, and absorb $R^\dagger$ into site $i-1$:
	\begin{equation}
		M^{[i-1]}_{\alpha_{i-2},\sigma_{i-1},\tilde\alpha_{i-1}}
		\leftarrow
		\sum_{\alpha_{i-1}}
		M^{[i-1]}_{\alpha_{i-2},\sigma_{i-1},\alpha_{i-1}}
		(R^\dagger)_{\alpha_{i-1},\tilde\alpha_{i-1}}.
	\end{equation}
	
	\section{SVD truncation}
	
	Given a matrix $X\in\mathbb{C}^{m\times n}$ with reduced SVD
	\begin{equation}
		X = U S V^\dagger,
	\end{equation}
	where
	\[
	U\in\mathbb{C}^{m\times r},\quad
	S=\mathrm{diag}(s_1,\dots,s_r),\quad
	V^\dagger\in\mathbb{C}^{r\times n},\quad
	s_1\ge s_2\ge \cdots \ge s_r\ge 0,
	\]
	truncate to bond dimension $D$ by keeping the $D$ largest singular values:
	\begin{equation}
		X \approx U_D S_D V_D^\dagger.
	\end{equation}
	
	The discarded weight is
	\begin{equation}
		\eta = \sum_{j>D} s_j^2.
	\end{equation}
	If the state is normalized, $\eta$ is the squared 2-norm of the discarded Schmidt tail.
	
	\begin{algorithm}[H]
		\caption{TruncateBySVD}
		\begin{algorithmic}[1]
			\Require matrix $X$, max bond dimension $D_{\max}$, tolerance $\eps$
			\Ensure truncated $U_D,S_D,V_D^\dagger$, discarded weight $\eta$
			\State compute reduced SVD: $X = U S V^\dagger$
			\State let $D = \min\{D_{\max},\#(s_j > \eps)\}$
			\If{$D=0$}
			\State $D\gets 1$
			\EndIf
			\State keep first $D$ singular values/vectors
			\State $\eta \gets \sum_{j>D} s_j^2$
		\end{algorithmic}
	\end{algorithm}
	
	\paragraph{Practical note.}
	In production codes, the retained bond dimension is often chosen using both a hard cap $D_{\max}$ and a discarded-weight tolerance, rather than only a singular-value threshold. A robust rule is to retain the smallest $D\le D_{\max}$ such that the discarded weight is below a prescribed tolerance, while always keeping at least one singular value.
	
	\section{Heisenberg MPO}
	
	For the spin-$1/2$ nearest-neighbor XXZ chain with field,
	\begin{equation}
		H =
		\sum_{i=1}^{L-1}
		\left[
		\frac{J}{2}\left(S_i^+ S_{i+1}^- + S_i^- S_{i+1}^+\right)
		+
		J_z S_i^z S_{i+1}^z
		\right]
		-
		h\sum_{i=1}^L S_i^z,
	\end{equation}
	a standard OBC MPO with bond dimension $\chi=5$ is
	\begin{equation}
		\hat W_{\mathrm{bulk}} =
		\begin{pmatrix}
			\Id & 0 & 0 & 0 & 0 \\
			S^+ & 0 & 0 & 0 & 0 \\
			S^- & 0 & 0 & 0 & 0 \\
			S^z & 0 & 0 & 0 & 0 \\
			-hS^z & \frac{J}{2}S^- & \frac{J}{2}S^+ & J_z S^z & \Id
		\end{pmatrix}.
	\end{equation}
	
	Here $\hat W_{\mathrm{bulk}}$ is an operator-valued matrix. The tensor components $W^{[i]}_{\beta_{i-1},\beta_i,\sigma,\sigma'}$ are obtained by expanding each operator entry in the local basis:
	\[
	(\hat O)_{\sigma,\sigma'} = \bra{\sigma}\hat O\ket{\sigma'}.
	\]
	
	For the local spin-$1/2$ basis $\{\ket{\uparrow},\ket{\downarrow}\}$, use
	\[
	S^+ = \begin{pmatrix}0&1\\0&0\end{pmatrix},\qquad
	S^- = \begin{pmatrix}0&0\\1&0\end{pmatrix},\qquad
	S^z = \frac12\begin{pmatrix}1&0\\0&-1\end{pmatrix}.
	\]
	
	For OBC, choose:
	\begin{equation}
		\hat W^{[1]} =
		\begin{pmatrix}
			-hS^z & \frac{J}{2}S^- & \frac{J}{2}S^+ & J_z S^z & \Id
		\end{pmatrix},
	\end{equation}
	\begin{equation}
		\hat W^{[L]} =
		\begin{pmatrix}
			\Id \\ S^+ \\ S^- \\ S^z \\ -hS^z
		\end{pmatrix},
	\end{equation}
	and use $\hat W_{\mathrm{bulk}}$ on all intermediate sites, with the chosen lower-triangular convention.
	
	In code, this means:
	\begin{itemize}
		\item the left boundary tensor has shape $(1,5,d,d)$ and is the row vector selecting the bottom row of the bulk operator-valued MPO,
		\item the right boundary tensor has shape $(5,1,d,d)$ and is the column vector selecting the leftmost column of the bulk operator-valued MPO.
	\end{itemize}
	
	\paragraph{Recommended unit test.}
	For $L=2,3,4$, contract the MPO to a dense Hamiltonian and compare entrywise against the Hamiltonian constructed directly by Kronecker products of the local operators.
	
	\section{AKLT Hamiltonian MPO}
	
	For the spin-1 AKLT chain, the Hamiltonian is
	\begin{equation}
		H_{\mathrm{AKLT}}
		=
		\sum_{i=1}^{L-1}
		\left[
		\mathbf S_i\cdot \mathbf S_{i+1}
		+
		\frac13\left(\mathbf S_i\cdot \mathbf S_{i+1}\right)^2
		\right].
		\label{eq:aklt-hamiltonian}
	\end{equation}
	This is the Hamiltonian whose exact ground state is the AKLT MPS given later in these notes.
	
	\subsection{Local spin-1 operators}
	
	In the local basis
	\[
	\{\ket{+1},\ket{0},\ket{-1}\},
	\]
	use the standard spin-1 matrices
	\begin{equation}
		S^z
		=
		\begin{pmatrix}
			1 & 0 & 0 \\
			0 & 0 & 0 \\
			0 & 0 & -1
		\end{pmatrix},
	\end{equation}
	\begin{equation}
		S^+
		=
		\sqrt{2}
		\begin{pmatrix}
			0 & 1 & 0 \\
			0 & 0 & 1 \\
			0 & 0 & 0
		\end{pmatrix},
		\qquad
		S^-
		=
		\sqrt{2}
		\begin{pmatrix}
			0 & 0 & 0 \\
			1 & 0 & 0 \\
			0 & 1 & 0
		\end{pmatrix}.
	\end{equation}
	
	Define the bilinear bond operator
	\begin{equation}
		X_{i,i+1}
		\equiv
		\mathbf S_i\cdot \mathbf S_{i+1}
		=
		\frac12\left(S_i^+S_{i+1}^- + S_i^-S_{i+1}^+\right)
		+
		S_i^z S_{i+1}^z.
		\label{eq:aklt-bilinear-bond}
	\end{equation}
	Then
	\begin{equation}
		h_{i,i+1}^{\mathrm{AKLT}}
		=
		X_{i,i+1}
		+
		\frac13 X_{i,i+1}^2,
		\label{eq:aklt-local-bond}
	\end{equation}
	and
	\[
	H_{\mathrm{AKLT}} = \sum_{i=1}^{L-1} h_{i,i+1}^{\mathrm{AKLT}}.
	\]
	
	\subsection{Factorized form of the bond interaction}
	
	Introduce the operator lists
	\begin{equation}
		O_1 = \frac{1}{\sqrt{2}}S^+,\qquad
		O_2 = \frac{1}{\sqrt{2}}S^-,\qquad
		O_3 = S^z,
	\end{equation}
	and
	\begin{equation}
		\bar O_1 = \frac{1}{\sqrt{2}}S^-,\qquad
		\bar O_2 = \frac{1}{\sqrt{2}}S^+,\qquad
		\bar O_3 = S^z.
	\end{equation}
	Then Eq.~\eqref{eq:aklt-bilinear-bond} becomes
	\begin{equation}
		X_{i,i+1}
		=
		\sum_{a=1}^3 O_a^{[i]} \bar O_a^{[i+1]}.
		\label{eq:X-factorized}
	\end{equation}
	Therefore
	\begin{equation}
		X_{i,i+1}^2
		=
		\sum_{a=1}^3 \sum_{b=1}^3
		\left(O_a O_b\right)^{[i]}
		\left(\bar O_a \bar O_b\right)^{[i+1]}.
		\label{eq:X2-factorized}
	\end{equation}
	Substituting Eqs.~\eqref{eq:X-factorized} and \eqref{eq:X2-factorized} into Eq.~\eqref{eq:aklt-local-bond}, we obtain
	\begin{equation}
		h_{i,i+1}^{\mathrm{AKLT}}
		=
		\sum_{a=1}^3 O_a^{[i]} \bar O_a^{[i+1]}
		+
		\frac13
		\sum_{a=1}^3 \sum_{b=1}^3
		\left(O_a O_b\right)^{[i]}
		\left(\bar O_a \bar O_b\right)^{[i+1]}.
		\label{eq:aklt-factorized-bond}
	\end{equation}
	
	\subsection{Exact MPO construction}
	
	Equation~\eqref{eq:aklt-factorized-bond} is an exact nearest-neighbor MPO. A convenient MPO has bond dimension
	\[
	\chi = 1 + 3 + 9 + 1 = 14.
	\]
	This slightly larger bond dimension is chosen because it leads to a completely explicit implementation without block-index ambiguities.
	
	\paragraph{Explicit MPO bond-index convention.}
	For implementation, the MPO bond index is ordered as follows:
	\begin{equation}
		\beta = 0
		\qquad \text{start channel},
	\end{equation}
	\begin{equation}
		\beta = 1,2,3
		\qquad \leftrightarrow \qquad
		a=1,2,3,
	\end{equation}
	\begin{equation}
		\beta = 4,\dots,12
		\qquad \leftrightarrow \qquad
		(a,b)\in
		\{(1,1),(1,2),(1,3),(2,1),(2,2),(2,3),(3,1),(3,2),(3,3)\},
		\label{eq:aklt-lexicographic-order}
	\end{equation}
	taken in lexicographic order, and
	\begin{equation}
		\beta = 13
		\qquad \text{terminal / identity-propagation channel}.
	\end{equation}
	This convention is used consistently in the bulk tensor and in both boundary tensors.
	
	Define
	\begin{equation}
		\bm O
		=
		\begin{pmatrix}
			O_1 & O_2 & O_3
		\end{pmatrix},
		\qquad
		\bar{\bm O}
		=
		\begin{pmatrix}
			\bar O_1 \\ \bar O_2 \\ \bar O_3
		\end{pmatrix},
	\end{equation}
	\begin{equation}
		\bm O^{(2)}
		=
		\begin{pmatrix}
			O_1 O_1 &
			O_1 O_2 &
			O_1 O_3 &
			O_2 O_1 &
			O_2 O_2 &
			O_2 O_3 &
			O_3 O_1 &
			O_3 O_2 &
			O_3 O_3
		\end{pmatrix},
	\end{equation}
	and
	\begin{equation}
		\bar{\bm O}^{(2)}
		=
		\begin{pmatrix}
			\bar O_1 \bar O_1 \\
			\bar O_1 \bar O_2 \\
			\bar O_1 \bar O_3 \\
			\bar O_2 \bar O_1 \\
			\bar O_2 \bar O_2 \\
			\bar O_2 \bar O_3 \\
			\bar O_3 \bar O_1 \\
			\bar O_3 \bar O_2 \\
			\bar O_3 \bar O_3
		\end{pmatrix}.
	\end{equation}
	The ordering of the $9$ entries in $\bm O^{(2)}$ and $\bar{\bm O}^{(2)}$ is the same lexicographic ordering defined in Eq.~\eqref{eq:aklt-lexicographic-order}.
	
	A bulk operator-valued MPO tensor is then
	\begin{equation}
		\hat W_{\mathrm{AKLT,bulk}}
		=
		\begin{pmatrix}
			\Id & \bm O & \frac13 \bm O^{(2)} & 0 \\
			0 & 0 & 0 & \bar{\bm O} \\
			0 & 0 & 0 & \bar{\bm O}^{(2)} \\
			0 & 0 & 0 & \Id
		\end{pmatrix},
		\label{eq:aklt-mpo-block}
	\end{equation}
	where the block sizes are $1|3|9|1$.
	
	In words:
	\begin{itemize}
		\item the channel $0\to a$ with $a=1,2,3$ starts the bilinear term via $O_a$,
		\item the channel $0\to(a,b)$ starts the biquadratic term via $\frac13 O_a O_b$,
		\item the channel $a\to 13$ finishes the bilinear term via $\bar O_a$,
		\item the channel $(a,b)\to 13$ finishes the biquadratic term via $\bar O_a \bar O_b$,
		\item the channel $13\to 13$ propagates the identity after a term has been completed.
	\end{itemize}
	With the index convention above, contracting the MPO generates exactly
	\[
	\sum_{i=1}^{L-1} h_{i,i+1}^{\mathrm{AKLT}}
	=
	\sum_{i=1}^{L-1}
	\left[
	\mathbf S_i\cdot \mathbf S_{i+1}
	+
	\frac13\left(\mathbf S_i\cdot \mathbf S_{i+1}\right)^2
	\right].
	\]
	
	For OBC, use boundary tensors
	\begin{equation}
		\hat W_{\mathrm{AKLT}}^{[1]}
		=
		\begin{pmatrix}
			\Id & \bm O & \frac13 \bm O^{(2)} & 0
		\end{pmatrix},
		\label{eq:aklt-left-boundary}
	\end{equation}
	and
	\begin{equation}
		\hat W_{\mathrm{AKLT}}^{[L]}
		=
		\begin{pmatrix}
			0 \\ \bar{\bm O} \\ \bar{\bm O}^{(2)} \\ \Id
		\end{pmatrix}.
		\label{eq:aklt-right-boundary}
	\end{equation}
	Thus, in the tensor convention of these notes,
	\begin{itemize}
		\item the left boundary tensor has shape $(1,14,3,3)$,
		\item each bulk tensor has shape $(14,14,3,3)$,
		\item the right boundary tensor has shape $(14,1,3,3)$.
	\end{itemize}
	The MPO physical indices are ordered as
	\[
	W^{[i]}_{\beta_{i-1},\beta_i,\sigma,\sigma'},
	\]
	with $\sigma$ the ket index and $\sigma'$ the bra index, exactly as in the Heisenberg MPO section.
	
	\subsection{Tensor components}
	
	As in the Heisenberg MPO section, the operator-valued entries are converted into numerical MPO tensor components by
	\begin{equation}
		W^{[i]}_{\beta_{i-1},\beta_i,\sigma,\sigma'}
		=
		\bra{\sigma}
		\hat W^{[i]}_{\beta_{i-1},\beta_i}
		\ket{\sigma'}.
	\end{equation}
	Here $\sigma,\sigma'\in\{+1,0,-1\}$ and the local physical dimension is $d=3$.
	
	\subsection{Implementation remarks}
	
	A transparent implementation strategy is:
	\begin{enumerate}
		\item construct the local spin-1 matrices $S^+,S^-,S^z$,
		\item define $O_a$ and $\bar O_a$,
		\item build the $9$ operators $O_a O_b$,
		\item build the $9$ operators $\bar O_a\bar O_b$,
		\item assemble the operator-valued boundary and bulk MPOs from
		Eqs.~\eqref{eq:aklt-mpo-block}, \eqref{eq:aklt-left-boundary}, and \eqref{eq:aklt-right-boundary},
		\item expand each operator entry into local matrix elements.
	\end{enumerate}
	
	\paragraph{Recommended unit tests.}
	For $L=2,3,4$:
	\begin{itemize}
		\item contract the MPO to a dense matrix,
		\item compare against the Hamiltonian built directly from Eq.~\eqref{eq:aklt-hamiltonian},
		\item verify that the AKLT MPS given later in these notes has energy
		\[
		E_0 = -\frac{2}{3}(L-1)
		\]
		for open chains, up to numerical precision.
	\end{itemize}
	
	\section{MPO environments}
	\label{sec:env}
	
	For a center at site $i$, let $L^{[i]}$ denote the contraction of sites
	$1,\dots,i-1$ into a left environment, and let $R^{[i]}$ denote the contraction
	of sites $i+1,\dots,L$ into a right environment.
	
	We store the left and right environments as tensors
	\[
	L^{[i]}_{\beta_{i-1}}(\alpha_{i-1},\alpha'_{i-1}),
	\qquad
	R^{[i]}_{\beta_i}(\alpha_i,\alpha'_i),
	\]
	with array shapes
	\[
	L^{[i]} \in \mathbb{C}^{\chi_{i-1}\times D_{i-1}\times D_{i-1}},
	\qquad
	R^{[i]} \in \mathbb{C}^{\chi_i\times D_i\times D_i}.
	\]
	
	\subsection{Left environment update}
	
	Suppose sites $1,\dots,i-1$ are represented by left-canonical tensors. Then the
	left environment is updated recursively by
	\begin{equation}
		L^{[i]}_{\beta_{i-1}}(\alpha_{i-1},\alpha'_{i-1})
		=
		\sum_{\substack{
				\alpha_{i-2},\alpha'_{i-2}\\
				\sigma,\sigma'\\
				\beta_{i-2}
		}}
		L^{[i-1]}_{\beta_{i-2}}(\alpha_{i-2},\alpha'_{i-2})
		A^{[i-1]}_{\alpha_{i-2},\sigma,\alpha_{i-1}}
		W^{[i-1]}_{\beta_{i-2},\beta_{i-1},\sigma,\sigma'}
		A^{[i-1]*}_{\alpha'_{i-2},\sigma',\alpha'_{i-1}}.
		\label{eq:left-env-consistent-rewritten}
	\end{equation}
	
	With array conventions
	\begin{itemize}
		\item \texttt{Lold[b,x,y]},
		\item \texttt{A[x,s,a]},
		\item \texttt{W[b,B,s,t]},
	\end{itemize}
	the validated NumPy implementation is
	\begin{verbatim}
		Lnew = np.einsum('bxy,xsa,bBst,ytc->Bac', Lold, A, W, A.conj())
	\end{verbatim}
	
	\subsection{Right environment update}
	
	Suppose sites $i+1,\dots,L$ are represented by right-canonical tensors. Then the
	right environment is updated recursively by
	\begin{equation}
		R^{[i]}_{\beta_i}(\alpha_i,\alpha'_i)
		=
		\sum_{\substack{
				\alpha_{i+1},\alpha'_{i+1}\\
				\sigma,\sigma'\\
				\beta_{i+1}
		}}
		B^{[i+1]}_{\alpha_i,\sigma,\alpha_{i+1}}
		W^{[i+1]}_{\beta_i,\beta_{i+1},\sigma,\sigma'}
		R^{[i+1]}_{\beta_{i+1}}(\alpha_{i+1},\alpha'_{i+1})
		B^{[i+1]*}_{\alpha'_i,\sigma',\alpha'_{i+1}}.
		\label{eq:right-env-consistent-rewritten}
	\end{equation}
	
	With array conventions
	\begin{itemize}
		\item \texttt{Rold[B,a,c]},
		\item \texttt{B[x,s,a]},
		\item \texttt{W[b,B,s,t]},
	\end{itemize}
	the validated NumPy implementation is
	\begin{verbatim}
		Rnew = np.einsum('xsa,bBst,Bac,ytc->bxy', B, W, Rold, B.conj())
	\end{verbatim}
	
	\subsection{Boundary initialization}
	
	For open boundary conditions,
	\begin{equation}
		L^{[1]}_{\beta_0}(1,1)=\delta_{\beta_0,1},
		\qquad
		R^{[L]}_{\beta_L}(1,1)=\delta_{\beta_L,1},
	\end{equation}
	with
	\[
	D_0=D_L=\chi_0=\chi_L=1.
	\]
	
	In code, these are simply
	\begin{verbatim}
		L = np.zeros((1,1,1), dtype=np.complex128); L[0,0,0] = 1
		R = np.zeros((1,1,1), dtype=np.complex128); R[0,0,0] = 1
	\end{verbatim}
	
	\subsection{Interpretation and convention warning}
	
	The update formulas above define the \emph{stored} left and right environment tensors.
	However, for the tensor, reshape, and MPO conventions used in these notes, the stored left
	environment does not enter the local effective-Hamiltonian action with the same bond-index
	order as it is stored.
	
	More precisely, if the stored left environment is represented in code as
	\[
	\texttt{L[b,x,y]},
	\]
	then in the validated one-site and two-site local contractions it must be used as
	\[
	\texttt{L[b,y,x]}.
	\]
	
	This point is easy to miss. A local operator built from environments can remain Hermitian
	even when this ordering is wrong, while still failing to agree with the exact projected
	operator
	\[
	P^\dagger H P.
	\]
	Therefore one should distinguish carefully between
	\begin{itemize}
		\item the recursion defining the stored environments, and
		\item the way those stored tensors are inserted into the local effective-Hamiltonian action.
	\end{itemize}
	
	\subsection{Practical recommendation}
	
	The left and right environment update routines should be tested independently against
	brute-force contractions on very small systems. After that, the resulting one-site and
	two-site environment-based local operators should be compared explicitly against dense
	projected references. In practice, this second test is the decisive one: it verifies not
	only that the environments are built recursively, but also that they are inserted into the
	local effective Hamiltonian with the correct index ordering.
	
	\section{One-site effective Hamiltonian}
	
	Let the center tensor at site $i$ be
	\[
	M^{[i]}_{\alpha_{i-1},\sigma_i,\alpha_i}.
	\]
	The local variational problem is
	\begin{equation}
		\mathcal{H}_{\mathrm{eff}}\,v = E v,
	\end{equation}
	where $v$ denotes the vectorized form of the center tensor.
	
	\paragraph{General structure.}
	The one-site effective Hamiltonian is obtained by contracting the full MPO with all
	MPS tensors outside site $i$, leaving only the local tensor $M^{[i]}$ open.
	In principle one may write this object symbolically as a contraction of a left environment,
	one local MPO tensor, and a right environment. In practice, however, the safest way to
	specify the operator---and the one that is directly useful for coding---is through its
	validated matrix-free contraction sequence.
	
	\paragraph{Convention-sensitive point.}
	For the conventions used in these notes, the stored left environment tensor does not enter
	the local action with the same MPS-bond index order as it is stored. If the stored array is
	\texttt{L[b,x,y]}, then the validated local contraction uses it as \texttt{L[b,y,x]}.
	Equivalently, the left environment enters the one-site effective Hamiltonian with its two
	MPS-bond indices interchanged relative to the most naive symbolic reading.
	
	\subsection{Validated matrix-free action}
	
	Let
	\[
	L^{[i]}_{\beta_{i-1}}(\alpha_{i-1},\alpha'_{i-1})
	\]
	be the stored left environment and
	\[
	R^{[i]}_{\beta_i}(\alpha_i,\alpha'_i)
	\]
	the stored right environment. Then the validated one-site matrix-free action is
	\begin{align}
		X_{\beta_{i-1},\alpha_{i-1},\sigma_i,\alpha_i'}
		&=
		\sum_{\alpha_{i-1}'}
		L^{[i]}_{\beta_{i-1}}(\alpha_{i-1}',\alpha_{i-1})
		M^{[i]}_{\alpha_{i-1}',\sigma_i,\alpha_i'},
		\label{eq:one-site-step1}
		\\
		Y_{\beta_i,\alpha_{i-1},\sigma_i',\alpha_i'}
		&=
		\sum_{\beta_{i-1},\sigma_i}
		W^{[i]}_{\beta_{i-1},\beta_i,\sigma_i',\sigma_i}
		X_{\beta_{i-1},\alpha_{i-1},\sigma_i,\alpha_i'},
		\label{eq:one-site-step2}
		\\
		(\mathcal{H}_{\mathrm{eff}}M)_{\alpha_{i-1},\sigma_i',\alpha_i}
		&=
		\sum_{\beta_i,\alpha_i'}
		Y_{\beta_i,\alpha_{i-1},\sigma_i',\alpha_i'}
		R^{[i]}_{\beta_i}(\alpha_i',\alpha_i).
		\label{eq:one-site-step3}
	\end{align}
	
	This is the convention-correct form for the one-site effective Hamiltonian in the present notes.
	In particular, the first step involves
	\[
	L^{[i]}_{\beta_{i-1}}(\alpha_{i-1}',\alpha_{i-1}),
	\]
	not
	\[
	L^{[i]}_{\beta_{i-1}}(\alpha_{i-1},\alpha_{i-1}').
	\]
	
	\subsection{Validated NumPy implementation}
	
	With arrays
	\begin{itemize}
		\item \texttt{L[b,x,y]} for the stored left environment,
		\item \texttt{M[y,s,z]} for the local center tensor,
		\item \texttt{W[b,B,s,t]} for the MPO tensor,
		\item \texttt{R[B,z,a]} for the stored right environment,
	\end{itemize}
	the validated implementation is
	\begin{verbatim}
		X  = np.einsum('byx,ysz->bxsz', L, M)
		Y  = np.einsum('bBst,bxsz->Bxtz', W, X)
		HM = np.einsum('Bxtz,Bza->xta', Y, R)
	\end{verbatim}
	
	Here \texttt{HM} has the same shape as \texttt{M}. The use of \texttt{'byx'} rather than
	\texttt{'bxy'} in the first contraction is essential.
	
	\paragraph{Interpretation of physical indices.}
	The MPO tensor is stored as
	\[
	W^{[i]}_{\beta_{i-1},\beta_i,\sigma,\sigma'},
	\]
	with the first physical index $\sigma$ the ket/output index and the second physical index
	$\sigma'$ the bra/input index. The contraction in Eq.~\eqref{eq:one-site-step2} is the
	validated routing consistent with this convention and with the environment storage described above.
	
	\subsection{Vectorization and local eigensolver interface}
	
	If
	\[
	M^{[i]} \in \mathbb{C}^{D_{i-1}\times d_i \times D_i},
	\]
	then the local vector-space dimension is
	\[
	N_{\mathrm{loc}} = D_{i-1} d_i D_i.
	\]
	The eigensolver vector
	\[
	v \in \mathbb{C}^{N_{\mathrm{loc}}}
	\]
	is identified with the tensor by
	\[
	M^{[i]} = \mathrm{reshape}(v,\,(D_{i-1},d_i,D_i)),
	\]
	using the fixed C-order convention of Section~\ref{sec:reshape}. After applying the local
	effective Hamiltonian, the output tensor is flattened back into a vector with the same convention.
	
	A practical wrapper is
	\begin{verbatim}
		from scipy.sparse.linalg import LinearOperator, eigsh
		
		def heff_one_site_matvec(v, L, W, R, Dl, d, Dr):
		M = v.reshape((Dl, d, Dr), order='C')
		X  = np.einsum('byx,ysz->bxsz', L, M)
		Y  = np.einsum('bBst,bxsz->Bxtz', W, X)
		HM = np.einsum('Bxtz,Bza->xta', Y, R)
		return HM.reshape(Dl * d * Dr, order='C')
		
		Nloc = Dl * d * Dr
		Heff = LinearOperator(
		shape=(Nloc, Nloc),
		matvec=lambda v: heff_one_site_matvec(v, L, W, R, Dl, d, Dr),
		dtype=np.complex128
		)
	\end{verbatim}
	
	The eigensolver should be initialized with the current local tensor:
	\begin{verbatim}
		v0 = M_current.reshape(Nloc, order='C')
	\end{verbatim}
	
	\subsection{Validation strategy}
	
	For small systems, one should explicitly build the dense projected operator
	\[
	P^\dagger H P
	\]
	for the one-site center and compare it against the dense matrix obtained by repeated
	application of the matrix-free routine above.
	
	\paragraph{Important warning.}
	Hermiticity of the environment-built local operator is necessary but not sufficient.
	With the conventions of these notes, a naive contraction using the stored left environment
	without the required index interchange can still produce a Hermitian operator, while failing
	to agree with the true projected operator \(P^\dagger H P\). Therefore the contraction
	sequence in Eqs.~\eqref{eq:one-site-step1}--\eqref{eq:one-site-step3} should be treated as
	the authoritative implementation for the present conventions.
	
	\section{Two-site effective Hamiltonian}
	
	At bond $(i,i+1)$, define the two-site center tensor
	\begin{equation}
		\Theta^{[i,i+1]}_{\alpha_{i-1},\sigma_i,\sigma_{i+1},\alpha_{i+1}},
		\label{eq:theta-def}
	\end{equation}
	so that
	\begin{equation}
		\ket{\psi}
		=
		\sum_{\alpha_{i-1},\sigma_i,\sigma_{i+1},\alpha_{i+1}}
		\Theta^{[i,i+1]}_{\alpha_{i-1},\sigma_i,\sigma_{i+1},\alpha_{i+1}}
		\ket{\alpha_{i-1}}_L \ket{\sigma_i}\ket{\sigma_{i+1}} \ket{\alpha_{i+1}}_R.
		\label{eq:two-site-state}
	\end{equation}
	
	The two-site local variational problem is the direct analogue of the standard
	two-site DMRG superblock problem, but expressed entirely in terms of MPS/MPO
	environments and a matrix-free local action.
	
	\paragraph{General structure.}
	In principle, the effective Hamiltonian at bond $(i,i+1)$ may be written symbolically
	as a contraction of
	\begin{itemize}
		\item a left environment built from sites $1,\dots,i-1$,
		\item the two local MPO tensors at sites $i$ and $i+1$,
		\item a right environment built from sites $i+2,\dots,L$.
	\end{itemize}
	In practice, however, the most reliable specification is again the validated contraction
	sequence itself.
	
	\paragraph{Convention-sensitive point.}
	For the conventions used in these notes, two implementation details are essential:
	\begin{enumerate}
		\item the stored left environment enters the local action with its two MPS-bond
		indices interchanged,
		\item the local MPO tensors must be contracted with the physical-index routing that
		is consistent with the validated implementation below.
	\end{enumerate}
	These points are easy to get wrong, and more naive contractions can produce Hermitian
	local operators that nevertheless fail to agree with the true projected operator
	\[
	P^\dagger H P.
	\]
	
	\subsection{Validated matrix-free action}
	
	Let
	\[
	L^{[i]}_{\beta_{i-1}}(\alpha_{i-1},\alpha'_{i-1})
	\]
	be the stored left environment built from sites $1,\dots,i-1$, and let
	\[
	R^{[i+1]}_{\beta_{i+1}}(\alpha_{i+1},\alpha'_{i+1})
	\]
	be the stored right environment built from sites $i+2,\dots,L$.
	
	The validated matrix-free two-site action is
	\begin{align}
		X_{\beta_{i-1},\alpha_{i-1},\sigma_i,\sigma_{i+1},\alpha_{i+1}'}
		&=
		\sum_{\alpha_{i-1}'}
		L^{[i]}_{\beta_{i-1}}(\alpha_{i-1}',\alpha_{i-1})
		\Theta_{\alpha_{i-1}',\sigma_i,\sigma_{i+1},\alpha_{i+1}'},
		\label{eq:two-site-step1}
		\\
		Y_{\beta_i,\alpha_{i-1},\sigma_i',\sigma_{i+1},\alpha_{i+1}'}
		&=
		\sum_{\beta_{i-1},\sigma_i}
		W^{[i]}_{\beta_{i-1},\beta_i,\sigma_i,\sigma_i'}
		X_{\beta_{i-1},\alpha_{i-1},\sigma_i,\sigma_{i+1},\alpha_{i+1}'},
		\label{eq:two-site-step2}
		\\
		Z_{\beta_{i+1},\alpha_{i-1},\sigma_i',\sigma_{i+1}',\alpha_{i+1}'}
		&=
		\sum_{\beta_i,\sigma_{i+1}}
		W^{[i+1]}_{\beta_i,\beta_{i+1},\sigma_{i+1},\sigma_{i+1}'}
		Y_{\beta_i,\alpha_{i-1},\sigma_i',\sigma_{i+1},\alpha_{i+1}'},
		\label{eq:two-site-step3}
		\\
		(\mathcal{H}_{\mathrm{eff}}\Theta)_{\alpha_{i-1},\sigma_i',\sigma_{i+1}',\alpha_{i+1}}
		&=
		\sum_{\beta_{i+1},\alpha_{i+1}'}
		Z_{\beta_{i+1},\alpha_{i-1},\sigma_i',\sigma_{i+1}',\alpha_{i+1}'}
		R^{[i+1]}_{\beta_{i+1}}(\alpha_{i+1}',\alpha_{i+1}).
		\label{eq:two-site-step4}
	\end{align}
	
	As in the one-site case, the first line uses
	\[
	L^{[i]}_{\beta_{i-1}}(\alpha_{i-1}',\alpha_{i-1})
	\]
	rather than
	\[
	L^{[i]}_{\beta_{i-1}}(\alpha_{i-1},\alpha_{i-1}').
	\]
	
	\subsection{Validated NumPy implementation}
	
	Let the two-site tensor be stored as
	\[
	\Theta[\alpha_{i-1},\sigma_i,\sigma_{i+1},\alpha_{i+1}],
	\]
	with shape
	\[
	(D_{i-1},d_i,d_{i+1},D_{i+1}).
	\]
	
	With arrays
	\begin{itemize}
		\item \texttt{L[b,x,y]} for the stored left environment,
		\item \texttt{Theta[y,u,v,z]} for the input two-site tensor,
		\item \texttt{W1[b,B,u,s]} for the MPO tensor on site $i$,
		\item \texttt{W2[B,C,v,t]} for the MPO tensor on site $i+1$,
		\item \texttt{R[C,z,a]} for the stored right environment,
	\end{itemize}
	the validated matrix-free contraction is
	\begin{verbatim}
		X  = np.einsum('byx,yuvz->bxuvz', L, Theta)
		Y  = np.einsum('bBus,bxuvz->Bxsvz', W1, X)
		Z  = np.einsum('BCvt,Bxsvz->Cxstz', W2, Y)
		HT = np.einsum('Cxstz,Cza->xsta', Z, R)
	\end{verbatim}
	
	The output tensor \texttt{HT[x,s,t,a]} has the same shape as \texttt{Theta}.
	
	\paragraph{Interpretation of physical indices.}
	The MPO tensors are stored as
	\[
	W^{[i]}_{\beta_{i-1},\beta_i,\sigma,\sigma'},
	\]
	with the first physical index the ket/output index and the second the bra/input index.
	The routing in the two-site contraction above is therefore convention-dependent and should
	be used exactly as written.
	
	\subsection{Vectorization and local eigensolver interface}
	
	If
	\[
	\Theta^{[i,i+1]} \in
	\mathbb{C}^{D_{i-1}\times d_i \times d_{i+1}\times D_{i+1}},
	\]
	then the local vector-space dimension is
	\[
	N_{\mathrm{loc}} = D_{i-1} d_i d_{i+1} D_{i+1}.
	\]
	The eigensolver vector
	\[
	v \in \mathbb{C}^{N_{\mathrm{loc}}}
	\]
	is identified with $\Theta$ by
	\[
	\Theta = \mathrm{reshape}(v,\,(D_{i-1},d_i,d_{i+1},D_{i+1})),
	\]
	again using the fixed C-order convention of Section~\ref{sec:reshape}.
	
	A practical wrapper is
	\begin{verbatim}
		def heff_two_site_matvec(v, L, W1, W2, R, Dl, d1, d2, Dr):
		Theta = v.reshape((Dl, d1, d2, Dr), order='C')
		X  = np.einsum('byx,yuvz->bxuvz', L, Theta)
		Y  = np.einsum('bBus,bxuvz->Bxsvz', W1, X)
		Z  = np.einsum('BCvt,Bxsvz->Cxstz', W2, Y)
		HT = np.einsum('Cxstz,Cza->xsta', Z, R)
		return HT.reshape(Dl * d1 * d2 * Dr, order='C')
	\end{verbatim}
	
	This routine is then passed to \texttt{LinearOperator} and solved by Lanczos or Davidson
	exactly as in the one-site case.
	
	\subsection{Validation strategy}
	
	For small systems, the dense matrix generated by repeated application of the matrix-free
	two-site routine above should be compared entrywise against the exact projected operator
	\[
	P^\dagger H P.
	\]
	This is the decisive correctness test for the environment-based implementation.
	
	\paragraph{Important warning.}
	Hermiticity alone is not a sufficient diagnostic. In the present conventions, one can
	obtain a two-site environment-built operator that is Hermitian and numerically stable,
	yet still incorrect because of an index-order mismatch in the left environment or in the
	MPO physical-index routing. Therefore the contraction sequence in
	Eqs.~\eqref{eq:two-site-step1}--\eqref{eq:two-site-step4} should be treated as the
	authoritative implementation for the present conventions.
	
	\section{Matrix-free eigensolver interface for local effective Hamiltonians}
	\label{sec:matfree}
	
	In practical DMRG implementations, the local effective Hamiltonian should \emph{never}
	be formed as a dense matrix except for unit tests on very small systems. Instead, one
	implements only the map
	\[
	v \mapsto \mathcal{H}_{\mathrm{eff}} v,
	\]
	and passes this map to an iterative eigensolver such as Lanczos or Davidson.
	
	\paragraph{Important convention warning.}
	For the tensor, MPO, and environment conventions adopted in these notes, the stored
	left environment tensor enters the local effective-Hamiltonian action with its two
	MPS-bond indices interchanged. In array form, if the stored left environment is
	\texttt{L[b,x,y]}, then the validated local matrix-free contractions use it as
	\texttt{L[b,y,x]}. This point is subtle but essential. A local operator built from
	environments can remain Hermitian even when this ordering is wrong, while still failing
	to agree with the true projected operator
	\[
	P^\dagger H P.
	\]
	Accordingly, the one-site and two-site matrix-free contractions given below should be
	treated as the authoritative implementations for the conventions of these notes.
	
	\subsection{One-site local solve}
	
	For a one-site center tensor
	\[
	M^{[i]} \in \mathbb{C}^{D_{i-1}\times d_i \times D_i},
	\]
	define the local vector-space dimension
	\[
	N_{\mathrm{loc}} = D_{i-1} d_i D_i.
	\]
	The eigensolver vector
	\[
	v\in\mathbb{C}^{N_{\mathrm{loc}}}
	\]
	is identified with the tensor $M^{[i]}$ by
	\begin{equation}
		M^{[i]} = \mathrm{reshape}(v,\,(D_{i-1},d_i,D_i)),
	\end{equation}
	using the fixed C-order convention of Section~\ref{sec:reshape}. After applying the
	local effective Hamiltonian, the result is reshaped back to a vector with the same convention.
	
	\paragraph{Validated one-site contraction.}
	The one-site matrix-free action is the validated contraction defined in
	Section~\ref{sec:one-site-heff}. In array form, with stored environment tensor
	\texttt{L[b,x,y]}, local tensor \texttt{M[y,s,z]}, MPO tensor \texttt{W[b,B,s,t]},
	and right environment \texttt{R[B,z,a]}, the authoritative contraction is
	\begin{verbatim}
		X  = np.einsum('byx,ysz->bxsz', L, M)
		Y  = np.einsum('bBst,bxsz->Bxtz', W, X)
		HM = np.einsum('Bxtz,Bza->xta', Y, R)
	\end{verbatim}
	
	\paragraph{Matrix-free wrapper.}
	A practical NumPy/\texttt{scipy} interface is:
	\begin{verbatim}
		from scipy.sparse.linalg import LinearOperator, eigsh
		
		def heff_one_site_matvec(v, L, W, R, Dl, d, Dr):
		M = v.reshape((Dl, d, Dr), order='C')
		X  = np.einsum('byx,ysz->bxsz', L, M)
		Y  = np.einsum('bBst,bxsz->Bxtz', W, X)
		HM = np.einsum('Bxtz,Bza->xta', Y, R)
		return HM.reshape(Dl * d * Dr, order='C')
		
		Nloc = Dl * d * Dr
		Heff = LinearOperator(
		shape=(Nloc, Nloc),
		matvec=lambda v: heff_one_site_matvec(v, L, W, R, Dl, d, Dr),
		dtype=np.complex128
		)
		
		E, vec = eigsh(Heff, k=1, which='SA', v0=v0)
		M = vec.reshape((Dl, d, Dr), order='C')
	\end{verbatim}
	
	Here \texttt{v0} should be the current local tensor flattened with the same reshape convention:
	\begin{verbatim}
		v0 = M_current.reshape(Nloc, order='C')
	\end{verbatim}
	
	\paragraph{Hermiticity and correctness checks.}
	For a Hermitian MPO and correctly built environments, the local operator should satisfy
	\[
	\langle x, \mathcal{H}_{\mathrm{eff}} y\rangle
	=
	\langle \mathcal{H}_{\mathrm{eff}} x, y\rangle
	\]
	up to numerical precision. However, Hermiticity alone is not a sufficient correctness test.
	For small systems, one should also compare the dense matrix obtained from the matrix-free
	routine against the exact projected operator
	\[
	P^\dagger H P.
	\]
	
	\subsection{Two-site local solve}
	
	For a two-site center tensor
	\[
	\Theta^{[i,i+1]} \in
	\mathbb{C}^{D_{i-1}\times d_i \times d_{i+1}\times D_{i+1}},
	\]
	the local vector-space dimension is
	\[
	N_{\mathrm{loc}} = D_{i-1} d_i d_{i+1} D_{i+1}.
	\]
	Again, the eigensolver vector $v$ is identified with the tensor $\Theta$ by
	\begin{equation}
		\Theta = \mathrm{reshape}(v,\,(D_{i-1},d_i,d_{i+1},D_{i+1})).
	\end{equation}
	After applying the local effective Hamiltonian, the output tensor is flattened back
	to a vector using the same C-order convention.
	
	\paragraph{Validated two-site contraction.}
	The two-site matrix-free action is the validated contraction defined in
	Section~\ref{sec:two-site-heff}. In array form, with stored left environment
	\texttt{L[b,x,y]}, two-site tensor \texttt{Theta[y,u,v,z]}, MPO tensors
	\texttt{W1[b,B,u,s]} and \texttt{W2[B,C,v,t]}, and stored right environment
	\texttt{R[C,z,a]}, the authoritative contraction is
	\begin{verbatim}
		X  = np.einsum('byx,yuvz->bxuvz', L, Theta)
		Y  = np.einsum('bBus,bxuvz->Bxsvz', W1, X)
		Z  = np.einsum('BCvt,Bxsvz->Cxstz', W2, Y)
		HT = np.einsum('Cxstz,Cza->xsta', Z, R)
	\end{verbatim}
	
	\paragraph{Matrix-free wrapper.}
	A practical wrapper is
	\begin{verbatim}
		def heff_two_site_matvec(v, L, W1, W2, R, Dl, d1, d2, Dr):
		Theta = v.reshape((Dl, d1, d2, Dr), order='C')
		X  = np.einsum('byx,yuvz->bxuvz', L, Theta)
		Y  = np.einsum('bBus,bxuvz->Bxsvz', W1, X)
		Z  = np.einsum('BCvt,Bxsvz->Cxstz', W2, Y)
		HT = np.einsum('Cxstz,Cza->xsta', Z, R)
		return HT.reshape(Dl * d1 * d2 * Dr, order='C')
	\end{verbatim}
	and this is passed to \texttt{LinearOperator} in the same way as for the one-site case.
	
	\paragraph{Practical remark.}
	The two-site contraction above was explicitly validated against the dense projected operator
	\[
	P^\dagger H P
	\]
	for small systems. This validation is essential, because more naive einsum routings can
	produce Hermitian local operators that are nevertheless incorrect.
	
	\subsection{Lanczos versus Davidson}
	
	For first implementations, Lanczos is usually sufficient and easiest to access through
	standard sparse-eigensolver interfaces. Davidson can be advantageous when the local
	dimension becomes larger or when a useful diagonal preconditioner is available.
	
	\begin{itemize}
		\item \textbf{Lanczos:} simplest robust choice for initial code.
		\item \textbf{Davidson:} often faster for large local problems, but requires more implementation effort.
	\end{itemize}
	
	In both cases, using the current local tensor as the initial guess is essential for sweep efficiency.
	
	\subsection{Dense reference for unit tests}
	
	For very small dimensions only, one should explicitly build the dense local projected operator
	\[
	H_{\mathrm{eff}} \in \mathbb{C}^{N_{\mathrm{loc}}\times N_{\mathrm{loc}}},
	\]
	and compare:
	\begin{itemize}
		\item dense matrix--vector multiplication versus the matrix-free routine,
		\item lowest eigenvalue from dense diagonalization versus Lanczos/Davidson,
		\item Hermiticity of the dense reference matrix,
		\item agreement of the environment-built operator with the exact projected operator
		\(P^\dagger H P\).
	\end{itemize}
	
	This unit test should be performed for both the one-site and two-site local solvers before
	running full sweeps.
	
	\section{Two-site DMRG}
	
	Two-site DMRG is the most robust finite-system ground-state algorithm and is recommended for a first implementation. It avoids many local-minimum problems of one-site DMRG and dynamically adjusts the bond dimension before truncation.
	
	\subsection{Two-site mixed-canonical form}
	
	At bond $(i,i+1)$, define left-orthonormal Schmidt states on sites $1,\dots,i-1$ and right-orthonormal Schmidt states on sites $i+2,\dots,L$:
	\[
	\ket{\alpha_{i-1}}_L,\qquad \ket{\alpha_{i+1}}_R.
	\]
	The variational two-site center tensor is
	\begin{equation}
		\Theta^{[i,i+1]}_{\alpha_{i-1},\sigma_i,\sigma_{i+1},\alpha_{i+1}},
		\label{eq:theta-def-alg}
	\end{equation}
	and the state is written as
	\begin{equation}
		\ket{\psi}
		=
		\sum_{\alpha_{i-1},\sigma_i,\sigma_{i+1},\alpha_{i+1}}
		\Theta^{[i,i+1]}_{\alpha_{i-1},\sigma_i,\sigma_{i+1},\alpha_{i+1}}
		\ket{\alpha_{i-1}}_L \ket{\sigma_i}\ket{\sigma_{i+1}} \ket{\alpha_{i+1}}_R.
		\label{eq:two-site-state-alg}
	\end{equation}
	
	This form is the direct analogue of the usual two-site DMRG superblock state and should be treated as the primary variational object. It is not necessary to define $\Theta$ through pre-existing tensors $A^{[i]}$ and $B^{[i+1]}$.
	
	\subsection{Solving the local problem}
	
	Given a matrix-free routine for $\Theta \mapsto \mathcal{H}_{\mathrm{eff}}\Theta$, solve for the lowest eigenpair using Lanczos or Davidson:
	\begin{equation}
		\mathcal{H}_{\mathrm{eff}} \Theta = E \Theta.
	\end{equation}
	
	The local vector space dimension is
	\begin{equation}
		N_{\mathrm{loc}} = D_{i-1} d_i d_{i+1} D_{i+1}.
	\end{equation}
	
	For spin-$1/2$ chains with uniform physical dimension $d_i=d_{i+1}=2$, this is
	\[
	N_{\mathrm{loc}} = 4 D_{i-1} D_{i+1}.
	\]
	
	The eigensolver should accept an initial guess. A good choice is the current two-site center tensor before optimization, flattened with the same C-order convention used everywhere else. This substantially accelerates convergence once sweeps are underway.
	
	\subsection{SVD split and truncation}
	
	After optimization, reshape the two-site tensor as a matrix:
	\begin{equation}
		\Theta_{(\alpha_{i-1}\sigma_i),(\sigma_{i+1}\alpha_{i+1})}
		\in \mathbb{C}^{(D_{i-1}d_i)\times(d_{i+1}D_{i+1})}.
		\label{eq:theta-matrix}
	\end{equation}
	
	Compute the reduced SVD:
	\begin{equation}
		\Theta = U S V^\dagger.
		\label{eq:theta-svd}
	\end{equation}
	
	Let
	\[
	D_{\mathrm{full}} = \mathrm{rank}(\Theta),\qquad
	D_{\mathrm{new}} = \min(D_{\max}, D_{\mathrm{full}})
	\]
	or more generally choose $D_{\mathrm{new}}$ from a singular-value threshold or discarded-weight tolerance. Keep only the largest $D_{\mathrm{new}}$ singular values.
	
	The discarded weight is
	\begin{equation}
		\eta^{[i]} = \sum_{j>D_{\mathrm{new}}} s_j^2.
		\label{eq:discarded-weight}
	\end{equation}
	
	For a left-to-right sweep, define
	\begin{equation}
		A^{[i]} \leftarrow \mathrm{reshape}\!\left(U,\,(D_{i-1},d_i,D_{\mathrm{new}})\right),
		\label{eq:two-site-left-update}
	\end{equation}
	\begin{equation}
		M^{[i+1]} \leftarrow \mathrm{reshape}\!\left(SV^\dagger,\,(D_{\mathrm{new}},d_{i+1},D_{i+1})\right).
		\label{eq:two-site-center-right}
	\end{equation}
	
	Then $A^{[i]}$ is left-canonical, and the orthogonality center has moved to site $i+1$.
	
	For a right-to-left sweep, instead use
	\begin{equation}
		M^{[i]} \leftarrow \mathrm{reshape}\!\left(US,\,(D_{i-1},d_i,D_{\mathrm{new}})\right),
		\label{eq:two-site-center-left}
	\end{equation}
	\begin{equation}
		B^{[i+1]} \leftarrow \mathrm{reshape}\!\left(V^\dagger,\,(D_{\mathrm{new}},d_{i+1},D_{i+1})\right),
		\label{eq:two-site-right-update}
	\end{equation}
	so that $B^{[i+1]}$ is right-canonical and the orthogonality center moves to site $i$.
	
	\paragraph{Important reshape note.}
	The reshapes in Eqs.~\eqref{eq:two-site-left-update}--\eqref{eq:two-site-right-update} are only correct if Eq.~\eqref{eq:theta-matrix} uses exactly the same flattening convention as the code.
	
	\subsection{Environment updates}
	
	After splitting the optimized two-site tensor, the environments must be updated so that
	the next local problem is built in the correct mixed-canonical gauge.
	
	\paragraph{Left-to-right sweep.}
	After obtaining the new left-canonical tensor $A^{[i]}$, update the left environment from
	site $i$ to site $i+1$:
	\begin{equation}
		L^{[i+1]}_{\beta_i}(\alpha_i,\alpha'_i)
		=
		\sum_{\substack{
				\alpha_{i-1},\alpha'_{i-1}\\
				\sigma_i,\sigma_i'\\
				\beta_{i-1}
		}}
		L^{[i]}_{\beta_{i-1}}(\alpha_{i-1},\alpha'_{i-1})
		A^{[i]}_{\alpha_{i-1},\sigma_i,\alpha_i}
		W^{[i]}_{\beta_{i-1},\beta_i,\sigma_i,\sigma_i'}
		A^{[i]*}_{\alpha'_{i-1},\sigma_i',\alpha'_i}.
	\end{equation}
	
	\paragraph{Right-to-left sweep.}
	After obtaining the new right-canonical tensor $B^{[i+1]}$, update the right environment
	from site $i+1$ to site $i$:
	\begin{equation}
		R^{[i]}_{\beta_i}(\alpha_i,\alpha'_i)
		=
		\sum_{\substack{
				\alpha_{i+1},\alpha'_{i+1}\\
				\sigma_{i+1},\sigma'_{i+1}\\
				\beta_{i+1}
		}}
		B^{[i+1]}_{\alpha_i,\sigma_{i+1},\alpha_{i+1}}
		W^{[i+1]}_{\beta_i,\beta_{i+1},\sigma_{i+1},\sigma'_{i+1}}
		R^{[i+1]}_{\beta_{i+1}}(\alpha_{i+1},\alpha'_{i+1})
		B^{[i+1]*}_{\alpha'_i,\sigma'_{i+1},\alpha'_{i+1}}.
	\end{equation}
	
	\paragraph{Convention reminder.}
	The formulas above define the \emph{stored} environments. As emphasized in
	Section~\ref{sec:env}, for the conventions of these notes the stored left environment enters
	the local effective-Hamiltonian action with its two MPS-bond indices interchanged relative
	to its storage order. Thus the environment-recursion formulas and the local-action formulas
	must not be conflated.
	
	\subsection{Full two-site sweep algorithm}
	
	\begin{algorithm}[H]
		\caption{TwoSiteDMRG}
		\begin{algorithmic}[1]
			\Require MPO $\{W^{[i]}\}$, initial MPS $\{M^{[i]}\}$, max bond dimension $D_{\max}$
			\Ensure approximate ground-state MPS and energy
			\State bring initial MPS into mixed-canonical form
			\State build all right environments $R^{[i]}$
			\State initialize left boundary environment $L^{[1]}$
			\Repeat
			\Comment{Left-to-right sweep}
			\For{$i=1$ to $L-1$}
			\State form the current two-site center tensor $\Theta^{[i,i+1]}$
			\State solve the lowest-eigenvalue problem $\mathcal{H}_{\mathrm{eff}}\Theta = E\Theta$
			\State reshape $\Theta$ as in Eq.~\eqref{eq:theta-matrix}
			\State compute reduced SVD $\Theta = U S V^\dagger$
			\State truncate to $D_{\mathrm{new}}\le D_{\max}$
			\State set $A^{[i]} \gets \mathrm{reshape}(U,(D_{i-1},d_i,D_{\mathrm{new}}))$
			\State set $M^{[i+1]} \gets \mathrm{reshape}(S V^\dagger,(D_{\mathrm{new}},d_{i+1},D_{i+1}))$
			\State record discarded weight $\eta^{[i]}$
			\State update left environment $L^{[i+1]}$
			\EndFor
			\Comment{Right-to-left sweep}
			\For{$i=L-1$ down to $1$}
			\State form the current two-site center tensor $\Theta^{[i,i+1]}$
			\State solve the lowest-eigenvalue problem
			\State reshape $\Theta$ as in Eq.~\eqref{eq:theta-matrix}
			\State compute reduced SVD $\Theta = U S V^\dagger$
			\State truncate to $D_{\mathrm{new}}\le D_{\max}$
			\State set $M^{[i]} \gets \mathrm{reshape}(US,(D_{i-1},d_i,D_{\mathrm{new}}))$
			\State set $B^{[i+1]} \gets \mathrm{reshape}(V^\dagger,(D_{\mathrm{new}},d_{i+1},D_{i+1}))$
			\State record discarded weight $\eta^{[i]}$
			\State update right environment $R^{[i]}$
			\EndFor
			\Until{energy, variance, and discarded-weight criteria are converged}
		\end{algorithmic}
	\end{algorithm}
	
	\subsection{Practical comments for coding}
	
	\begin{itemize}
		\item \textbf{This should be the first DMRG implementation.} It is much easier to debug than one-site DMRG.
		\item \textbf{Keep the orthogonality center explicit.} During a left-to-right sweep, after the split the center lives on site $i+1$ as $M^{[i+1]}$; during a right-to-left sweep, it lives on site $i$ as $M^{[i]}$.
		\item \textbf{Do not compare tensors directly across sweeps without gauge fixing.}
		\item \textbf{Always unit test the local effective-Hamiltonian matvec} against a dense reference for small dimensions.
		\item \textbf{Monitor the maximum discarded weight}
		\[
		\eta_{\max} = \max_i \eta^{[i]}
		\]
		during each sweep.
		\item \textbf{Benchmark first against exact diagonalization} for small Heisenberg chains before attempting larger systems.
	\end{itemize}
	
	\section{One-site DMRG}
	\label{sec:one-site-dmrg}
	
	One-site DMRG is best treated as an advanced refinement of the two-site algorithm. In practice,
	a robust implementation typically uses several two-site sweeps first and only then switches to
	one-site sweeps, optionally with subspace expansion or density-matrix perturbation.
	
	\subsection{One-site mixed-canonical form}
	
	In one-site DMRG, the MPS is kept in mixed-canonical form with center at site $i$:
	\begin{equation}
		\ket{\psi}
		=
		\sum_{\alpha_{i-1},\sigma_i,\alpha_i}
		M^{[i]}_{\alpha_{i-1},\sigma_i,\alpha_i}
		\ket{\alpha_{i-1}}_L \ket{\sigma_i} \ket{\alpha_i}_R.
	\end{equation}
	
	The local optimization is then performed over the center tensor $M^{[i]}$ while the left
	and right Schmidt bases are kept fixed.
	
	\subsection{One-site sweep algorithm}
	
	\begin{algorithm}[H]
		\caption{OneSiteDMRG}
		\begin{algorithmic}[1]
			\Require MPO $\{W^{[i]}\}$, initial MPS $\{M^{[i]}\}$, max bond dimension $D_{\max}$
			\Ensure approximate ground-state MPS and energy
			\State bring the initial MPS into mixed-canonical form
			\State build all right environments
			\State initialize the left boundary environment
			\Repeat
			\Comment{Left-to-right sweep}
			\For{$i=1$ to $L-1$}
			\State solve the one-site effective-Hamiltonian problem using the validated matrix-free contraction of Section~\ref{sec:matfree}
			\State reshape the optimized center tensor as a left-grouped matrix
			\State compute an SVD: $M^{[i]} = U S V^\dagger$
			\State retain at most $D_{\max}$ singular values
			\State reshape $U \to A^{[i]}$ as a left-canonical tensor
			\State absorb $S V^\dagger$ into $M^{[i+1]}$
			\State update $L^{[i+1]}$
			\EndFor
			\Comment{Right-to-left sweep}
			\For{$i=L$ down to $2$}
			\State solve the one-site effective-Hamiltonian problem using the validated matrix-free contraction of Section~\ref{sec:matfree}
			\State reshape the optimized center tensor as a right-grouped matrix
			\State compute an SVD: $M^{[i]} = U S V^\dagger$
			\State retain at most $D_{\max}$ singular values
			\State reshape $V^\dagger$ using the inverse of Eq.~\eqref{eq:right-group} to obtain a right-canonical tensor $B^{[i]}$
			\State absorb $U S$ into $M^{[i-1]}$
			\State update $R^{[i-1]}$
			\EndFor
			\Until{energy and variance are converged}
		\end{algorithmic}
	\end{algorithm}
	
	\paragraph{Remark.}
	In standard one-site DMRG, the SVD step mainly regauges the state and shifts the orthogonality
	center; unlike two-site DMRG, it does not naturally enlarge the accessible bond space. This is
	one reason why one-site DMRG is more prone to local minima unless combined with subspace expansion
	or density-matrix perturbation.
	
	\paragraph{Convention reminder.}
	For the conventions of these notes, the one-site local matrix-free action used inside one-site
	DMRG must be the validated contraction summarized in Section~\ref{sec:matfree}, rather than a
	naive contraction built directly from the stored left environment without the required bond-index
	interchange.
	
	\paragraph{Practical warning.}
	One-site DMRG is prone to local minima and should not be used as a first standalone implementation.
	A robust implementation should either:
	\begin{itemize}
		\item start with several two-site sweeps, or
		\item use a subspace-expansion / density-matrix perturbation method.
	\end{itemize}
	
	\section{AKLT benchmark state}
	An exact MPS representation of the AKLT state is, up to gauge, given by
	\begin{align}
		A^{+} &= \begin{pmatrix} 0 & \sqrt{2/3} \\ 0 & 0 \end{pmatrix},\\
		A^{0} &= \begin{pmatrix} -1/\sqrt{3} & 0 \\ 0 & 1/\sqrt{3} \end{pmatrix},\\
		A^{-} &= \begin{pmatrix} 0 & 0 \\ -\sqrt{2/3} & 0 \end{pmatrix}.
	\end{align}
	For this gauge, the tensors satisfy
	\begin{equation}
		\sum_{\sigma\in\{+,0,-\}} A^\sigma A^{\sigma\dagger} = \Id_2,
	\end{equation}
	so they are \emph{right-canonical} in the convention of Eq.~\eqref{eq:right-canonical}.
	This state is an excellent benchmark:
	\begin{itemize}
		\item the AKLT ground state admits an exact MPS representation with bond dimension $D=2$,
		\item for open boundary conditions, the AKLT Hamiltonian has a fourfold degenerate ground-state manifold associated with spin-$1/2$ edge modes, and every ground state in that manifold has exact energy
		\[
		E_0 = -\frac23 (L-1).
		\]
	\end{itemize}
	
	\section{Convergence diagnostics}
	
	Reliable convergence diagnostics for DMRG should include:
	\begin{itemize}
		\item energy change
		\[
		|\Delta E| < \eps_E,
		\]
		\item energy variance
		\[
		\sigma_H^2 = \bra{\psi}H^2\ket{\psi} - \bra{\psi}H\ket{\psi}^2,
		\]
		\item maximum discarded weight during a sweep,
		\item stability of entanglement spectrum or local observables.
	\end{itemize}
	
	\paragraph{Practical note.}
	For small systems, the variance can be benchmarked against dense linear algebra. In production calculations, it can be obtained either from an MPO representation of $H^2$ or by first computing $H\ket{\psi}$ as an MPS and then evaluating its norm. One should keep in mind that an MPO for $H^2$ can have a significantly larger bond dimension than the MPO for $H$.
	
	\paragraph{Important caution.}
	Direct tensor-by-tensor comparison is not gauge-invariant and should \emph{not} be used as a primary convergence criterion unless the gauge is fixed identically before comparison.
	
	\section{Implementation checklist}
	
	Before production calculations, the implementation should be validated in a strictly
	controlled order. The purpose of this checklist is not only to test individual routines,
	but also to ensure that the tensor-index conventions, reshape conventions, and environment
	contractions are mutually consistent.
	
	\begin{enumerate}
		\item \textbf{Canonicalization test.}
		Starting from random tensors or a random MPS, perform left- and right-canonicalization
		sweeps and verify the defining relations
		\begin{align}
			\sum_{\alpha_{i-1},\sigma_i}
			A^{[i]}_{\alpha_{i-1},\sigma_i,\alpha_i}
			A^{[i]*}_{\alpha_{i-1},\sigma_i,\alpha_i'}
			&= \delta_{\alpha_i,\alpha_i'},
			\\
			\sum_{\sigma_i,\alpha_i}
			B^{[i]}_{\alpha_{i-1},\sigma_i,\alpha_i}
			B^{[i]*}_{\alpha_{i-1}',\sigma_i,\alpha_i}
			&= \delta_{\alpha_{i-1},\alpha_{i-1}'}.
		\end{align}
		These tests verify both the QR/SVD logic and the fixed reshape conventions.
		
		\item \textbf{MPO construction test.}
		For small chains, such as $L=2,3,4$, contract the MPO to a dense matrix and compare
		it entrywise against a Hamiltonian built directly from Kronecker products of local operators.
		This test should be performed for both the Heisenberg MPO and the AKLT MPO.
		
		\item \textbf{Environment-recursion test.}
		Verify that the left and right environment updates reproduce brute-force contractions
		on very small systems. This confirms that the recursive environment-building routines are
		internally consistent.
		
		\item \textbf{Hermiticity test of local matrix-free operators.}
		For random small tensors, verify numerically that
		\[
		\langle x,\mathcal{H}_{\mathrm{eff}} y\rangle
		=
		\langle \mathcal{H}_{\mathrm{eff}} x, y\rangle
		\]
		up to numerical precision, for both one-site and two-site matrix-free local actions.
		
		\item \textbf{One-site projected-operator test.}
		For a small mixed-canonical MPS, explicitly construct the dense projected one-site operator
		\[
		P^\dagger H P
		\]
		and compare it against the dense matrix obtained by repeated application of the matrix-free
		one-site environment-based local action.
		
		For the conventions of these notes, the validated one-site contraction is
		\begin{verbatim}
			X  = np.einsum('byx,ysz->bxsz', L, M)
			Y  = np.einsum('bBst,bxsz->Bxtz', W, X)
			HM = np.einsum('Bxtz,Bza->xta', Y, R)
		\end{verbatim}
		
		\item \textbf{Two-site projected-operator test.}
		For a small mixed-canonical MPS, explicitly construct the dense projected two-site operator
		\[
		P^\dagger H P
		\]
		and compare it against the dense matrix obtained by repeated application of the matrix-free
		two-site environment-based local action.
		
		For the conventions of these notes, the validated two-site contraction is
		\begin{verbatim}
			X  = np.einsum('byx,yuvz->bxuvz', L, Theta)
			Y  = np.einsum('bBus,bxuvz->Bxsvz', W1, X)
			Z  = np.einsum('BCvt,Bxsvz->Cxstz', W2, Y)
			HT = np.einsum('Cxstz,Cza->xsta', Z, R)
		\end{verbatim}
		
		This comparison against the dense projected operator is the decisive local correctness test.
		
		\item \textbf{Two-site DMRG benchmark for Heisenberg chains.}
		Run finite-system two-site DMRG for small spin-$1/2$ Heisenberg chains and compare the
		converged MPS energy against exact diagonalization. Agreement to numerical precision on
		small systems is the main global benchmark for the implementation.
		
		\item \textbf{AKLT benchmark.}
		Validate the AKLT implementation in two ways:
		\begin{itemize}
			\item compare the DMRG ground-state energy against the exact open-chain result
			\[
			E_0 = -\frac23 (L-1),
			\]
			\item evaluate the energy of the exact AKLT MPS and verify that it reproduces the same result.
		\end{itemize}
		
		\item \textbf{Convergence diagnostics in production runs.}
		After the implementation is validated locally and globally, monitor in larger calculations:
		\begin{itemize}
			\item energy change between sweeps,
			\item maximum discarded weight,
			\item eventual energy variance,
			\item stability of local observables and entanglement spectra.
		\end{itemize}
	\end{enumerate}
	
	\paragraph{Most important practical warning.}
	For the conventions used in these notes, Hermiticity of the environment-built local operator
	is \emph{not} sufficient to guarantee correctness. In particular, an incorrect treatment of the
	stored left environment or an incorrect MPO physical-index routing can still yield a Hermitian
	local operator, yet one that differs from the true projected operator \(P^\dagger H P\).
	Therefore no full DMRG sweeps should be trusted until the one-site and two-site local operators
	have been compared explicitly against dense projected references on small systems.
	
	\paragraph{Recommended debugging order.}
	If a full DMRG calculation produces stable but incorrect energies, the safest debugging sequence is:
	\begin{enumerate}
		\item verify canonicalization,
		\item verify MPO-to-dense agreement,
		\item verify left and right environment updates separately,
		\item verify the one-site local projected operator,
		\item verify the two-site local projected operator,
		\item only then test full sweeps.
	\end{enumerate}
	In practice, the projected-operator tests are usually more informative than Hermiticity tests alone.
	
	\paragraph{Authoritative implementation templates.}
	For the conventions of these notes, the one-site and two-site matrix-free contraction templates
	written above should be treated as authoritative reference implementations. Future optimizations,
	refactorings, or changes of contraction order should always be benchmarked back against these
	forms on small systems before being used in production calculations.
	
	\section{Practical coding recommendation}
	
	For a first reliable implementation, the following strategy is strongly recommended.
	
	\begin{itemize}
		\item \textbf{Implement two-site DMRG first.}
		Among standard finite-system algorithms, two-site DMRG is the most robust place to begin.
		It is substantially easier to debug than one-site DMRG, naturally allows bond dimensions to
		adapt before truncation, and provides a direct setting in which the local variational space is
		large enough to avoid many of the local-minimum problems of one-site updates.
		
		\item \textbf{Fix tensor conventions once and use them everywhere.}
		In particular, fix:
		\begin{itemize}
			\item the MPS tensor convention
			\[
			A^{[i]}_{\alpha_{i-1},\sigma_i,\alpha_i},
			\]
			\item the MPO tensor convention
			\[
			W^{[i]}_{\beta_{i-1},\beta_i,\sigma_i,\sigma_i'},
			\]
			with the first physical index the ket/output index and the second the bra/input index,
			\item the reshape rules used in QR, SVD, vectorization, and inverse reshaping,
			\item the memory order, e.g.\ NumPy C-order.
		\end{itemize}
		Small inconsistencies in these choices are among the most common sources of silent errors.
		
		\item \textbf{Use reduced QR for regauging and SVD for truncation.}
		Reduced QR without pivoting is convenient for moving orthogonality centers and constructing
		canonical forms. SVD should be used when a truncation decision is required, because it gives
		direct access to Schmidt singular values and discarded weights.
		
		\item \textbf{Keep the orthogonality-center structure explicit.}
		In two-site DMRG, the variational object is the two-site tensor
		\[
		\Theta^{[i,i+1]}_{\alpha_{i-1},\sigma_i,\sigma_{i+1},\alpha_{i+1}},
		\]
		and the tensors to the left and right should be kept in the corresponding mixed-canonical form.
		During a left-to-right sweep, after splitting the optimized two-site tensor, the center moves
		to the right; during a right-to-left sweep, it moves to the left. This structure should be
		reflected explicitly in the code.
		
		\item \textbf{Do not build dense local effective Hamiltonians in production sweeps.}
		Except for very small validation tests, the local effective Hamiltonian should be accessed only
		through a matrix-free map
		\[
		v \mapsto \mathcal{H}_{\mathrm{eff}} v.
		\]
		This is the only scalable strategy and is the natural interface to Lanczos or Davidson.
		
		\item \textbf{Treat the validated contraction templates as authoritative.}
		For the conventions of these notes, the one-site and two-site local environment-based
		contractions are subtle. In particular, the stored left environment enters the local action
		with its MPS-bond indices interchanged, and the MPO physical-index routing in the two-site
		action must be implemented exactly as in the validated contraction sequence.
		These convention-dependent details should not be re-derived from memory during coding.
		
		\item \textbf{Benchmark against dense projected operators before full sweeps.}
		Before running finite-system DMRG sweeps, one should verify on small systems that the
		environment-built one-site and two-site local operators agree with the exact projected operators
		\[
		P^\dagger H P.
		\]
		This is a more stringent and more informative test than Hermiticity alone.
		
		\item \textbf{Benchmark against exact diagonalization on small systems.}
		After the local-operator tests pass, run two-site DMRG for small Heisenberg chains and compare
		the converged energy against exact diagonalization. This is the standard global correctness test.
		
		\item \textbf{Use the AKLT chain as an exact MPS benchmark.}
		The AKLT model is particularly valuable because:
		\begin{itemize}
			\item the Hamiltonian has an exact MPO representation of modest bond dimension,
			\item the ground state admits an exact MPS representation with bond dimension $D=2$,
			\item the open-chain ground-state energy is known exactly,
			\[
			E_0 = -\frac23 (L-1).
			\]
		\end{itemize}
		Therefore the AKLT chain tests simultaneously the MPO, the local solver, the sweep logic,
		and the handling of exact low-bond-dimension MPS ground states.
		
		\item \textbf{Use the current local tensor as the eigensolver initial guess.}
		In finite-system sweeps, the previously optimized local tensor provides an excellent initial
		vector for Lanczos or Davidson. This is a major practical improvement in convergence speed.
		
		\item \textbf{Monitor discarded weights and, when available, variances.}
		The energy alone is not always a sufficient convergence indicator. In production calculations,
		one should monitor at least:
		\begin{itemize}
			\item the energy change between sweeps,
			\item the maximum discarded weight during each sweep,
			\item the eventual energy variance,
			\item stability of local observables and entanglement spectra.
		\end{itemize}
	\end{itemize}
	
	\paragraph{Recommended order of development.}
	A robust development sequence is:
	\begin{enumerate}
		\item local operators,
		\item MPO construction,
		\item canonicalization routines,
		\item environment updates,
		\item one-site projected-operator validation,
		\item two-site projected-operator validation,
		\item two-site DMRG benchmark for Heisenberg chains,
		\item AKLT benchmark,
		\item only afterward: one-site DMRG, variance evaluation, and further optimizations.
	\end{enumerate}
	
	\paragraph{Most important practical warning.}
	For the conventions used in these notes, a local operator built from environments can be
	Hermitian and numerically stable while still being incorrect. This can happen if the stored
	left environment is inserted into the local action with the wrong bond-index order, or if
	the MPO physical-index routing in the local contraction is inconsistent with the chosen MPO
	storage convention. Therefore full DMRG sweeps should \emph{not} be trusted until the
	matrix-free local operator has been compared explicitly against the dense projected operator
	\(P^\dagger H P\) on small systems.
	
	\paragraph{Final recommendation.}
	In practical coding work, the safest philosophy is:
	\begin{quote}
		first validate the local operator, then validate small-system energies, and only then scale up.
	\end{quote}
	For environment-based DMRG implementations, this order is not merely convenient; it is essential.
	
	\end{document}